% !TEX root = ../elteikthesis_en.tex
\chapter{Background and Related Work}
\label{ch:background}

This chapter surveys the technical foundations and prior art that motivate the design decisions made in this project. Section~\ref{sec:bg-llm} introduces large language models and their limitations for domain-specific question answering. Section~\ref{sec:bg-small} discusses small and locally deployed models and the role of quantization. Section~\ref{sec:bg-rag} explains retrieval-augmented generation and why it is preferred over fine-tuning for this use case. Section~\ref{sec:bg-related} reviews existing university chatbot deployments. Section~\ref{sec:bg-tech} justifies the specific technology choices made for this project.

% ─────────────────────────────────────────────────────────────
\section{Large Language Models}
\label{sec:bg-llm}

Large language models (LLMs) are neural networks with billions of parameters trained on vast text corpora to predict the next token in a sequence~\cite{brown2020gpt3}. The dominant architecture is the Transformer~\cite{vaswani2017attention}, which uses multi-head self-attention to relate every token in a sequence to every other token, allowing the model to capture long-range dependencies that earlier recurrent architectures struggled with.

During pre-training, a model such as GPT or Llama sees hundreds of billions of tokens drawn from web pages, books, and code repositories. This process implicitly encodes a large amount of world knowledge into the model's weights. At inference time the model can answer questions, summarise text, translate languages, and generate code without any task-specific training.

\subsection{Limitations for Domain-Specific Question Answering}

Despite their broad capabilities, general-purpose LLMs exhibit several weaknesses that are directly relevant to the ELTE chatbot use case.

\paragraph{Knowledge cutoff.} Pre-training terminates at a fixed point in time; events and documents published after that date are unknown to the model. For a faculty information system that changes each academic semester, a static knowledge cutoff is unacceptable.

\paragraph{Hallucination.} LLMs are trained to produce fluent, plausible text, not necessarily factually correct text. When the model lacks reliable knowledge about a topic it will often fabricate plausible-sounding but incorrect facts~\cite{lewis2020rag}. In an academic advising context a hallucinated prerequisite or registration deadline could directly harm students.

\paragraph{Context window limits.} Although recent models accept increasingly long inputs, the context window is still finite. A model cannot be given an entire document corpus in its prompt; some form of selection is required.

\paragraph{Privacy and data sovereignty.} Accessing a cloud-hosted LLM API means sending potentially sensitive queries about students and courses to a third-party server. Many institutional environments — including the current one — prohibit this under data protection policies.

% ─────────────────────────────────────────────────────────────
\section{Small and Local Language Models}
\label{sec:bg-small}

One response to the privacy and cost concerns of cloud APIs is to run a smaller model locally. The Llama family of models~\cite{touvron2023llama2}, developed by Meta AI and released under a permissive licence, covers a range of sizes from 1B to 70B parameters. Smaller variants (3B--8B parameters) can run on consumer hardware within an 8~GB RAM budget if their weights are quantised.

\subsection{Quantization}

A full-precision (FP32) 3B-parameter model requires roughly 12~GB of memory — more than is available on most laptops. \emph{Quantization} reduces memory by representing each weight with fewer bits. The GGUF file format~\cite{gerganov2023llama}, used by the llama.cpp inference engine and popularised by Ollama, supports 4-bit and 8-bit quantization schemes. A 4-bit quantised Llama~3.2~3B model fits comfortably within 2--3~GB of RAM, leaving headroom for the operating system and embedding model.

The trade-off is a small but measurable degradation in output quality. Empirical benchmarks show that 4-bit GGUF quantization reduces scores on standard reasoning benchmarks by 2--5\% relative to the full-precision baseline~\cite{dettmers2023qlora}. For a document-grounded Q\&A system where the LLM is asked to summarise retrieved text rather than perform novel reasoning, this degradation is acceptable.

\subsection{Ollama}

Ollama is an open-source application that wraps llama.cpp behind a simple REST API, manages model downloads, and handles hardware detection (CPU-only vs.\ GPU-offloaded inference). It exposes a single endpoint \texttt{POST /api/generate} that accepts a prompt and streams back a response. This makes it straightforward to integrate into a Python backend without managing low-level inference libraries directly.

% ─────────────────────────────────────────────────────────────
\section{Retrieval-Augmented Generation}
\label{sec:bg-rag}

Retrieval-augmented generation (RAG) was formalised by Lewis et al.~\cite{lewis2020rag} as a method that conditions an LLM's output on documents retrieved at query time. Instead of relying solely on knowledge encoded in model weights, a RAG pipeline retrieves a small number of relevant passages from an external corpus and inserts them directly into the prompt, giving the model grounded evidence to draw on.

\subsection{RAG versus Fine-Tuning}

An alternative to RAG is to fine-tune the LLM on domain documents, updating its weights to encode institutional knowledge. Table~\ref{tab:rag-vs-ft} compares the two approaches for this use case.

\begin{table}[H]
\centering
\caption{RAG versus fine-tuning for a document-grounded chatbot}
\label{tab:rag-vs-ft}
\begin{tabular}{lll}
\hline
\textbf{Property} & \textbf{RAG} & \textbf{Fine-Tuning} \\
\hline
Knowledge update & Re-ingest new documents & Retrain the model \\
Compute required & CPU inference only & GPU, hours to days \\
Hallucination risk & Low (grounded in retrieved text) & Moderate \\
Source attribution & Built-in (chunk metadata) & Not straightforward \\
Out-of-scope refusal & Possible (no relevant chunks found) & Difficult to enforce \\
\hline
\end{tabular}
\end{table}

For a university information system whose documents change each semester, the ability to update the knowledge base by simply re-ingesting new files — without touching the model weights — is decisive. Fine-tuning a 3B-parameter model also requires a GPU that is not available in the target deployment environment. RAG is therefore the preferred architecture.

\subsection{Embedding-Based Retrieval}

The retrieval step in a RAG pipeline works as follows. During indexing, every text chunk is encoded into a dense vector (an \emph{embedding}) by a bi-encoder model. At query time the user's question is encoded with the same model, and the chunks whose embeddings are most similar to the query embedding are selected. Cosine similarity or approximate nearest-neighbour search is used to rank candidates efficiently.

The embedding model used in this project is \texttt{all-MiniLM-L6-v2}~\cite{reimers2019sbert}, a 22M-parameter Sentence-Transformer that maps text to 384-dimensional vectors. It was chosen for three reasons: it is small enough to run on CPU without perceptible latency, it consistently ranks among the top models on the MTEB retrieval benchmarks for its size class~\cite{muennighoff2022mteb}, and it is natively supported by the \texttt{sentence-transformers} Python library.

Retrieved chunks are stored in ChromaDB, a lightweight open-source vector database that runs as an in-process library with no separate server process. ChromaDB persists its index to disk, supports metadata filtering, and requires no configuration beyond pointing it at a directory — properties that suit a single-machine deployment.

\subsection{Prompt Augmentation}

After retrieval the RAG backend constructs a prompt that contains the retrieved chunks and the user's question. The LLM is instructed to answer only from the provided context and to state when it cannot find an answer. This instruction, combined with the low temperature setting (0.1), keeps the model's output factual and tightly grounded.

% ─────────────────────────────────────────────────────────────
\section{Related Work: University Chatbots}
\label{sec:bg-related}

AI-powered chatbots are being adopted in higher education at an accelerating pace. A 2025 survey found that 37\% of colleges now provide institution-wide chatbot licences and 57\% of higher education leaders regard AI as a strategic priority~\cite{edtech2025chatbots}. At the same time, 14\% of institutions have built homegrown systems, often to address concerns that commercial products cannot satisfy — chief among them data privacy and vendor lock-in.

\subsection{University of Michigan: UM-GPT}

The University of Michigan launched UM-GPT in August 2023, making it one of the first large universities to deploy a campus-wide generative AI platform at scale. The infrastructure runs in a private Microsoft Azure environment with on-premises options for sensitive workloads. To reduce hallucination, the platform incorporates retrieval-augmented generation over institutional document repositories, mirroring the core approach used in this thesis. The university also deployed complementary tools: Maizey, a codeless RAG builder that has produced over 3{,}500 department-level chatbot instances, and GoToCollege, a scholarship information assistant~\cite{edtech2025chatbots}.

\subsection{University of California, Irvine: ZotGPT}

UC Irvine launched ZotGPT in January 2024. Within its first year the platform attracted 15{,}000 unique users and expanded to offer access to multiple LLMs, an API tier for researchers, and ZotGPT ClassChat for course-specific knowledge bases. The infrastructure spans both Azure AI and Amazon Web Services. Chris Price, the university's AI architect, noted that institutional deployment substantially reduced barriers for stakeholders who were previously reluctant to experiment with AI tools~\cite{edtech2025chatbots}.

\subsection{Arizona State University: CreateAI}

Arizona State University's CreateAI platform, launched in October 2023, aggregates over 50 generative AI models under a single interface and provides a MyAI Builder tool for custom deployments. Specialist bots have been built for behavioural health training and philosophy education. The platform is hosted on AWS. Elizabeth Reilley, ASU's AI acceleration director, emphasised that well-designed guidelines foster rather than inhibit innovation~\cite{edtech2025chatbots}.

\subsection{Positioning This Work}

All three deployments above share two characteristics: they rely on cloud infrastructure (Azure or AWS) and they target large, research-intensive universities with dedicated IT teams and substantial budgets. The system developed in this thesis occupies a different point in the design space: it runs entirely offline on a single laptop with no cloud dependencies, uses a quantised 3B-parameter model rather than a commercial API, and is designed for a single faculty rather than a whole university. This makes it feasible to deploy in environments where cloud data transfer is prohibited and where no dedicated IT infrastructure is available.

% ─────────────────────────────────────────────────────────────
\section{Technology Choices}
\label{sec:bg-tech}

\subsection{Language Model: Llama~3.2~3B via Ollama}

Llama~3.2~3B was selected as the generation model because it fits within the 8~GB RAM constraint after 4-bit quantization, it is freely available under Meta's community licence, and it achieves competitive performance on instruction-following benchmarks for its size class. Ollama provides a stable HTTP interface that decouples the FastAPI backend from the specifics of the inference engine; switching to a different model requires only changing a single configuration value.

\subsection{Embedding Model: all-MiniLM-L6-v2}

The \texttt{all-MiniLM-L6-v2} model encodes both documents and queries into 384-dimensional vectors. Its small size (22M parameters, $\approx$90~MB on disk) means that embedding a query adds only a few milliseconds to request latency on CPU. Critically, the same model must be used during both index construction and query time; a mismatch would produce incompatible vector spaces and entirely wrong retrieval results.

\subsection{Vector Database: ChromaDB}

ChromaDB runs as a Python library with no network daemon, stores its index as files on disk, and exposes a simple collection API. For a single-node deployment with a corpus of approximately 6{,}000 chunks, its performance is more than adequate and its operational simplicity is a significant advantage over alternatives such as Milvus or Weaviate that require separate server processes.

\subsection{Backend Framework: FastAPI}

FastAPI is a modern Python web framework built on Starlette and Pydantic. It provides automatic OpenAPI documentation, asynchronous request handling, and type-checked request and response models. These properties make it straightforward to expose the RAG pipeline as a REST API that can be consumed by both the browser-based chat UI and command-line clients.
